% GENERATED FILE. Edit canonical lessons, then run npm run generate:books.
% canonical-source-hash: fnv1a64:3284efe5357aa415
% canonical-lessons: AR-C493-nadima, AR-C493-nadim, AR-C493-shafaqa, AR-C493-ataba, AR-C493-yas

\chapter{Regret and pity}
\label{ch:ar-regret-and-pity}
\hlchaptermodality{493}

\begin{chapteropening}
\textbf{By the end of this chapter:} \emph{I can say that someone regretted something or gently reproached someone, and talk about pity and despair.}

Complete the last lesson of chapter 493: I can say that someone regretted something or gently reproached someone, and talk about pity and despair.
\end{chapteropening}

\section[\texorpdfstring{\ar{ندم} (nadima)}{ندم (nadima)}]{\ar{ندم} (nadima) — to regret}

\label{lesson:AR-C493-nadima}



\begin{quote}
Before the new one: say the Arabic for to envy, then the Arabic for to be jealous.
\end{quote}

\subsection*{You'll want to know: \ar{ندم}}
\textbf{\ar{ندم}} — \emph{nadima} — \textquotedblleft{}to regret\textquotedblright{}.

\begin{quote}
\ar{ندم} \ar{الولد} \ar{لأنه} \ar{كذب} \ar{على} \ar{أمه}.
\end{quote}

\begin{quote}
\emph{nadima al-walad li-annahu kadhaba ʿalā ummihi.}
\end{quote}

\begin{quote}
The boy was sorry because he had lied to his mother.
\end{quote}

Present: \ar{يندم}, \emph{yandam}.

It takes \ar{على}: \ar{ندم} \ar{على} \ar{ما} \ar{قال}.

\begin{cousinweb}
\textbf{\ar{ندم}} (\emph{nadima}) is built on the native Arabic root \emph{n-d-m}, \textquotedblleft{}to repent, to regret\textquotedblright{} (Lane). The noun nadam is regret.
\end{cousinweb}

\begin{grammarlens}[title={What you've built}]
One more word for describing what you lived and felt.
\end{grammarlens}

\subsection*{Your turn}
\begin{itemize}
\raggedright
  \item \emph{Say it:} \emph{nadima}
  \item \emph{Say it:} \emph{nadima}, once more
  \item \emph{Say it:} \emph{nadima al-walad li-annahu kadhaba ʿalā ummihi.} — The boy was sorry because he had lied to his mother.
  \item \emph{Recall:} say the Arabic for to envy, then the Arabic for to be jealous, then say \emph{nadima} again
\end{itemize}

\begin{tcolorbox}[breakable,colback=teal!4,colframe=teal!35!black,title={Before you move on}]
What does \ar{ندم} mean? (\textquotedblleft{}To regret\textquotedblright{}.) Say it once more.
\end{tcolorbox}
\section[\texorpdfstring{\ar{نادم} (nādim)}{نادم (nādim)}]{\ar{نادم} (nādim) — regretful}

\label{lesson:AR-C493-nadim}



\begin{quote}
Before the new one: say the Arabic for to be jealous, then the Arabic for to regret.
\end{quote}

\subsection*{You'll want to know: \ar{نادم}}
\textbf{\ar{نادم}} — \emph{nādim} — \textquotedblleft{}regretful\textquotedblright{}.

\begin{quote}
\ar{كان} \ar{اللاعب} \ar{نادما} \ar{بعد} \ar{المباراة}.
\end{quote}

\begin{quote}
\emph{kāna al-lāʿib nādiman baʿda al-mubārā.}
\end{quote}

\begin{quote}
The player was full of regret after the match.
\end{quote}

\begin{cousinweb}
\textbf{\ar{نادم}} (\emph{nādim}) is built on the native Arabic root \emph{n-d-m}, \textquotedblleft{}to repent, to regret\textquotedblright{} (Lane). Nādim is the active participle: one who regrets.
\end{cousinweb}

\begin{grammarlens}[title={What you've built}]
One more word for describing what you lived and felt.
\end{grammarlens}

\subsection*{Your turn}
\begin{itemize}
\raggedright
  \item \emph{Say it:} \emph{nādim}
  \item \emph{Say it:} \emph{nādim}, once more
  \item \emph{Say it:} \emph{kāna al-lāʿib nādiman baʿda al-mubārā.} — The player was full of regret after the match.
  \item \emph{Recall:} say the Arabic for to be jealous, then the Arabic for to regret, then say \emph{nādim} again
\end{itemize}

\begin{tcolorbox}[breakable,colback=teal!4,colframe=teal!35!black,title={Before you move on}]
What does \ar{نادم} mean? (\textquotedblleft{}Regretful\textquotedblright{}.) Say it once more.
\end{tcolorbox}
\section[\texorpdfstring{\ar{شفقة} (shafaqa)}{شفقة (shafaqa)}]{\ar{شفقة} (shafaqa) — pity}

\label{lesson:AR-C493-shafaqa}



\begin{quote}
Before the new one: say the Arabic for to regret, then the Arabic for regretful.
\end{quote}

\subsection*{You'll want to know: \ar{شفقة}}
\textbf{\ar{شفقة}} — \emph{shafaqa} — \textquotedblleft{}pity\textquotedblright{}.

\begin{quote}
\ar{شعرت} \ar{بالشفقة} \ar{على} \ar{الحصان} \ar{المتعب}.
\end{quote}

\begin{quote}
\emph{shaʿartu bi-sh-shafaqa ʿalā al-ḥiṣān al-mutʿab.}
\end{quote}

\begin{quote}
I felt pity for the tired horse.
\end{quote}

It takes \ar{على} for the one pitied.

\begin{cousinweb}
\textbf{\ar{شفقة}} (\emph{shafaqa}) is built on the native Arabic root \emph{sh-f-q}, \textquotedblleft{}to be anxious for, to pity\textquotedblright{} (Lane). Shafaqa is a tender worry for someone.
\end{cousinweb}

\begin{grammarlens}[title={What you've built}]
One more word for describing what you lived and felt.
\end{grammarlens}

\subsection*{Your turn}
\begin{itemize}
\raggedright
  \item \emph{Say it:} \emph{shafaqa}
  \item \emph{Say it:} \emph{shafaqa}, once more
  \item \emph{Say it:} \emph{shaʿartu bi-sh-shafaqa ʿalā al-ḥiṣān al-mutʿab.} — I felt pity for the tired horse.
  \item \emph{Recall:} say the Arabic for to regret, then the Arabic for regretful, then say \emph{shafaqa} again
\end{itemize}

\begin{tcolorbox}[breakable,colback=teal!4,colframe=teal!35!black,title={Before you move on}]
What does \ar{شفقة} mean? (\textquotedblleft{}Pity\textquotedblright{}.) Say it once more.
\end{tcolorbox}
\section[\texorpdfstring{\ar{عاتب} (ʿātaba)}{عاتب (ʿātaba)}]{\ar{عاتب} (ʿātaba) — to reproach (gently)}

\label{lesson:AR-C493-ataba}



\begin{quote}
Before the new one: say the Arabic for regretful, then the Arabic for pity.
\end{quote}

\subsection*{You'll want to know: \ar{عاتب}}
\textbf{\ar{عاتب}} — \emph{ʿātaba} — \textquotedblleft{}to reproach (gently)\textquotedblright{}.

\begin{quote}
\ar{عاتبت} \ar{الأم} \ar{ابنها} \ar{لأنه} \ar{لم} \ar{يتصل} \ar{بها}.
\end{quote}

\begin{quote}
\emph{ʿātabat al-umm ibnahā li-annahu lam yattaṣil bihā.}
\end{quote}

\begin{quote}
The mother reproached her son because he had not phoned her.
\end{quote}

Present: \ar{يعاتب}, \emph{yuʿātib}.

It is a gentle, loving complaint, not anger.

\begin{cousinweb}
\textbf{\ar{عاتب}} (\emph{ʿātaba}) is built on the native Arabic root \emph{ʿ-t-b}, \textquotedblleft{}to blame, to be angry with\textquotedblright{} (Lane). Lane gives yaʿtibu for Form I; Form III ʿātaba is to reproach someone face to face.
\end{cousinweb}

\begin{grammarlens}[title={What you've built}]
One more word for describing what you lived and felt.
\end{grammarlens}

\subsection*{Your turn}
\begin{itemize}
\raggedright
  \item \emph{Say it:} \emph{ʿātaba}
  \item \emph{Say it:} \emph{ʿātaba}, once more
  \item \emph{Say it:} \emph{ʿātabat al-umm ibnahā li-annahu lam yattaṣil bihā.} — The mother reproached her son because he had not phoned her.
  \item \emph{Recall:} say the Arabic for regretful, then the Arabic for pity, then say \emph{ʿātaba} again
\end{itemize}

\begin{tcolorbox}[breakable,colback=teal!4,colframe=teal!35!black,title={Before you move on}]
What does \ar{عاتب} mean? (\textquotedblleft{}To reproach (gently)\textquotedblright{}.) Say it once more.
\end{tcolorbox}
\section[\texorpdfstring{\ar{يأس} (yaʾs)}{يأس (yaʾs)}]{\ar{يأس} (yaʾs) — despair}

\label{lesson:AR-C493-yas}



\begin{quote}
Before the new one: say the Arabic for pity, then the Arabic for to reproach.
\end{quote}

\subsection*{You'll want to know: \ar{يأس}}
\textbf{\ar{يأس}} — \emph{yaʾs} — \textquotedblleft{}despair\textquotedblright{}.

\begin{quote}
\ar{لم} \ar{يشعر} \ar{المريض} \ar{باليأس} \ar{أبدا}.
\end{quote}

\begin{quote}
\emph{lam yashʿur al-marīḍ bi-l-yaʾs abadan.}
\end{quote}

\begin{quote}
The patient never felt despair.
\end{quote}

\begin{cousinweb}
\textbf{\ar{يأس}} (\emph{yaʾs}) is built on the native Arabic root \emph{y-ʾ-s}, \textquotedblleft{}to despair, to give up hope\textquotedblright{} (Lane). Yaʾs is the noun of giving up hope.
\end{cousinweb}

\begin{grammarlens}[title={What you've built}]
One more word for describing what you lived and felt.
\end{grammarlens}

\subsection*{Your turn}
\begin{itemize}
\raggedright
  \item \emph{Say it:} \emph{yaʾs}
  \item \emph{Say it:} \emph{yaʾs}, once more
  \item \emph{Say it:} \emph{lam yashʿur al-marīḍ bi-l-yaʾs abadan.} — The patient never felt despair.
  \item \emph{Recall:} say the Arabic for pity, then the Arabic for to reproach, then say \emph{yaʾs} again
\end{itemize}

\begin{tcolorbox}[breakable,colback=teal!4,colframe=teal!35!black,title={Before you move on}]
What does \ar{يأس} mean? (\textquotedblleft{}Despair\textquotedblright{}.) Say it once more.
\end{tcolorbox}
